\section*{الفصل \ref{ch:g5:weather-and-climate} --- الطقس والمناخ}

\begin{solution}{exo:g5:weather-and-climate:1}
(أ) \omterm{def:g1:watching-the-weather:weather}{طقس} (هذا الصباح فقط). (ب) \omterm{def:g5:weather-and-climate:climate}{مناخ} (كل صيف). (ج) \omterm{def:g1:watching-the-weather:weather}{طقس}
(بعد ظهر أمس). (د) \omterm{def:g5:weather-and-climate:climate}{مناخ} (كل شهر، سنة بعد سنة).
\end{solution}

\begin{solution}{exo:g5:weather-and-climate:2}
\omterm{def:g5:weather-and-climate:climate}{مناخ} مكان ما هو \omterm{def:g1:watching-the-weather:weather}{الطقس} الذي يعرفه عادة، شهرًا بعد شهر، على مدى
سنين كثيرة. ويأخذ علماء \omterm{def:g1:watching-the-weather:weather}{الطقس} متوسط قياسات 30 سنة.
\end{solution}

\begin{solution}{exo:g5:weather-and-climate:3}
أبرد الأشهر ديسمبر (نحو \babelsublr{6~$^\circ$C}؛ ويناير يكاد يكون بالبرودة
نفسها). وأكثرها مطرًا نوفمبر (نحو 160~mm). لا: المنحنى لا ينزل أبدًا
تحت \babelsublr{0~$^\circ$C}.
\end{solution}

\begin{solution}{exo:g5:weather-and-climate:4}
يوليو أدفأ الأشهر، بنحو \babelsublr{5~$^\circ$C}. ولا يوجد شهر أدفأ من
\babelsublr{10~$^\circ$C}.
\end{solution}

\begin{solution}{exo:g5:weather-and-climate:5}
\omterm{def:g3:continents-and-oceans:equator}{خط الاستواء}: النطاق الحار. \omterm{def:g3:continents-and-oceans:poles}{القطب الشمالي}: النطاق البارد (الشمالي
منهما). ومنتصف المسافة بين \omterm{def:g3:continents-and-oceans:equator}{خط الاستواء} \omterm{def:g3:continents-and-oceans:poles}{والقطب الجنوبي}: النطاق
المعتدل الجنوبي (فهو يقع بين \omterm{def:g5:weather-and-climate:tropics}{مدار} الجدي \omterm{def:g5:weather-and-climate:tropics}{والدائرة القطبية}
الجنوبية).
\end{solution}

\begin{solution}{exo:g5:weather-and-climate:6}
هيلو: النطاق الحار. فينيكس وسياتل: النطاق المعتدل الشمالي.
أوتكياغفيك: النطاق البارد الشمالي، وراء \omterm{def:g5:weather-and-climate:tropics}{الدائرة القطبية} الشمالية.
\end{solution}

\begin{solution}{exo:g5:weather-and-climate:7}
السنة كلها: $147 + 96 + 106 + 81 + 48 + 37 + 15 + 25 + 41 + 99 + 160 + 145
= 1\,000$ مليمتر. وديسمبر ويناير وفبراير:
$145 + 147 + 96 = 388$ مليمترًا، أي أكثر من ثلث \omterm{def:g4:clouds-rain-snow:precipitation}{مطر} السنة في
ثلاثة أشهر.
\end{solution}

\begin{solution}{exo:g5:weather-and-climate:8}
فينيكس: $35 - 14 = 21$ درجة. هيلو: $25 - 22 = 3$ درجات. فالفصول
تغيّر درجة الحرارة أكثر بكثير في فينيكس؛ أما في هيلو فتكاد تكون هي نفسها
طوال السنة.
\end{solution}

\begin{solution}{exo:g5:weather-and-climate:9}
\omterm{def:g5:weather-and-climate:climate}{مناخ} \omterm{def:g5:weather-and-climate:rainforest}{الغابة المطيرة}: فالجو حار طوال السنة (27 أو \babelsublr{28~$^\circ$C}، بلا
شتاء)، ويسقط \omterm{def:g4:clouds-rain-snow:precipitation}{مطر} كثير في كل شهر (160~mm على الأقل، وأكثر
من 2\,000~mm في السنة). والأرجح أن المدينة في النطاق الحار، قرب
\omterm{def:g3:continents-and-oceans:equator}{خط الاستواء}.
\end{solution}

\begin{solution}{exo:g5:weather-and-climate:10}
التفسير الأول: الشاطئ في النطاق الحار، حيث الجو دافئ طوال السنة،
والحديقة المثلجة في نطاق معتدل أو بارد، حيث يقع يناير في
الشتاء. والتفسير الثاني: الشاطئ في النصف الجنوبي من الأرض،
حيث يقع يناير في الصيف، والحديقة المثلجة في النصف الشمالي،
حيث يقع يناير في الشتاء: فالفصول معكوسة بين نصفي
الكرة الأرضية.
\end{solution}

\begin{solution}{exo:g5:weather-and-climate:11}
درجات الحرارة: فينيكس أحر بكثير، نحو \babelsublr{35~$^\circ$C} في يوليو مقابل
نحو \babelsublr{20~$^\circ$C} في أغسطس لسياتل، ويناير فيها (نحو
\babelsublr{14~$^\circ$C}) دافئ بقدر مايو في سياتل. \omterm{def:g4:clouds-rain-snow:precipitation}{والمطر}: أقل من 200~mm في السنة
في فينيكس، وهي \omterm{def:g5:weather-and-climate:desert}{صحراء}، مقابل نحو 1\,000~mm في سياتل. فليس لهما
\omterm{def:g5:weather-and-climate:climate}{المناخ} نفسه. والنطاقات تخبرنا، تقريبًا، كم تسخّن الشمس
مكانًا ما، من \omterm{def:g3:continents-and-oceans:equator}{خط الاستواء} إلى \omterm{def:g3:continents-and-oceans:poles}{القطبين}؛ لكنها لا تخبرنا كم يسقط من
\omterm{def:g4:clouds-rain-snow:precipitation}{المطر}، ولا بأثر \omterm{def:g2:water-on-earth:ocean}{البحر}، أو \omterm{def:g1:landscapes-around-us:landform}{الجبال}، أو الرياح.
\end{solution}
